\documentclass[11pt]{amsart}
\usepackage{amsmath,amssymb,amsthm}
\usepackage[margin=1.1in]{geometry}
\usepackage{array,longtable,booktabs}
\usepackage{hyperref}
\emergencystretch=2em

\newtheorem{theorem}{Theorem}
\newtheorem{lemma}[theorem]{Lemma}
\newtheorem{proposition}[theorem]{Proposition}
\theoremstyle{remark}
\newtheorem{remark}[theorem]{Remark}

\newcommand{\Inf}{\mathrm{Inf}}
\newcommand{\Var}{\mathrm{Var}}
\newcommand{\E}{\mathbb{E}}

\title{Erasure entropies of sparse linear DNFs}
\author{Yitzchak Shmalo}
\address{Einstein Institute of Mathematics, The Hebrew University of Jerusalem, Jerusalem, Israel}
\date{}

\begin{document}

\begin{abstract}
For a Boolean function $f$ and a coordinate $i$, the erasure entropy $E_i(f)$ is the conditional entropy of the event
$i\in U$ given $U\setminus\{i\}$, where $U$ is drawn from the spectral distribution $\hat f^2$. The Fourier
entropy-influence conjecture is equivalent to every non-constant $f$ having a coordinate with $E_i(f)=O(\Inf_i(f))$. We
consider monotone formulas in disjunctive normal form (DNFs) whose terms all have width $w$, any two of which share at most one variable, and in which every
variable lies in at most $D$ terms. We show that $E_i(f)/\Inf_i(f)=2+o(1)$ at every relevant coordinate, uniformly in
the number of variables and terms, as $w\to\infty$ with $D\le2^{(1-\eta)(w-1)}$ for a fixed $\eta>0$, and we give an
explicit error bound for each $w$, $D$ and degree of $i$. The proof rests on a separation inequality valid for every
monotone function: if the derivative $D_if$ puts little spectral weight outside a family $\Omega$ of sets and the
restriction $f|_{x_i=0}$ puts little weight on $\Omega$, then $E_i/\Inf_i$ is close to $2$. For a linear DNF, $\Omega$
is the family of sets that contain many of the other variables of some term through $i$.
\end{abstract}

\subjclass[2020]{06E30, 42C10, 94A17}
\keywords{Fourier entropy-influence conjecture, Boolean functions, Fourier entropy, total influence}

\maketitle

\section{Introduction}\label{sec:intro}

For $f\colon\{0,1\}^n\to\{-1,1\}$ write $f=\sum_S\hat f(S)\chi_S$ with $\chi_S(x)=(-1)^{\sum_{j\in S}x_j}$, and let $U$
be a random subset with $\Pr[U=S]=\hat f(S)^2$. The Fourier entropy is $H(f)=H(U)$ and the total influence is
$I(f)=\E|U|=\sum_i\Inf_i(f)$, where $\Inf_i(f)=\Pr[i\in U]$. The Fourier entropy-influence (FEI) conjecture of Friedgut
and Kalai~\cite{FK96} asks for an absolute $C$ with $H(f)\le C\,I(f)$ for every Boolean $f$. It is known for symmetric
functions~\cite{OWZ} and for read-once formulas, over gates of bounded arity~\cite{OT13} and over AND and OR
gates~\cite{CKLS}. It was proved for read-$k$ decision trees with $k$ fixed in~\cite{WWW} and for linear threshold
functions in~\cite{ShmLTF}. Its min-entropy
version is known for read-$k$ DNFs with $k$ fixed~\cite{Shalev,ACKSW}, and Han~\cite{Han} bounds $H(f)$ by a multiple of
$I(f)$ plus $\sum_i\Inf_i(f)\log_2(1/\Inf_i(f))$, with the multiple lowered to the sharp value $1/\ln 2$ by Li
and Han~\cite{LH25}. Kelman, Kindler, Lifshitz, Minzer and Safra~\cite{KKLMS} prove a
weaker form of the conjecture: the spectrum of a Boolean function of constant variance is concentrated on
$2^{O(I(f)\log I(f))}$ characters.

The erasure entropy (in the sense of Verd\'u and Weissman~\cite{VW08}) of $f$ at $i$ is
\[
E_i(f)=H\bigl(\mathbf 1[i\in U]\bigm| U\setminus\{i\}\bigr).
\]
It is the part of the entropy drop from $f$ to its two restrictions at $x_i$ that is not controlled by a constant times
$\Inf_i$~\cite[Remarks~3.1--3.2]{Han},~\cite[(32)]{GMP}, and the FEI conjecture holds if and only if there is an
absolute $c$ such that every non-constant Boolean function has a coordinate with $E_i(f)\le c\,\Inf_i(f)$~\cite{Shm1}.
The trivial bound $E_i\le h_2(\Inf_i)$, where $h_2(p)=-p\log_2p-(1-p)\log_2(1-p)$ is the binary entropy, loses a factor $\log(1/\Inf_i)$; the question is when that factor can be
removed, and at which coordinates.

This paper computes the ratio $E_i/\Inf_i$ asymptotically for a class of DNFs in which the coordinates are linked by
many terms. Let $F$ be a monotone DNF on $\{0,1\}^n$. We call $F$ \emph{linear} if any two of its terms share at most
one variable, \emph{$w$-uniform} if all its terms have $w$ variables, and of \emph{degree at most $D$} if every
variable lies in at most $D$ terms. Put $f=1-2F$. Read-once DNFs are the linear DNFs of degree one; a linear DNF of
degree $D$ is read-$D$, but its terms may overlap in a pattern with many cycles, for instance the $m$ rows and $m$
columns of an $m\times m$ array.

\begin{theorem}\label{thm:main}
Fix $\eta\in(0,1)$. Over all $n$, all $w$-uniform linear monotone DNFs $F$ on $\{0,1\}^n$ of degree at most
$D\le2^{(1-\eta)(w-1)}$, and all coordinates $i$ with $\Inf_i(f)>0$,
\[
\sup\Bigl|\frac{E_i(f)}{\Inf_i(f)}-2\Bigr|\longrightarrow0\qquad(w\to\infty).
\]
In particular $\sum_iE_i(f)=(2+o(1))\,I(f)$ on this class.
\end{theorem}

Theorem~\ref{thm:main} follows from an explicit inequality, Theorem~\ref{thm:lin}, valid for each width, degree and
coordinate. The value $2$ is the limit of $E_i/\Inf_i$ for a coordinate inside one long AND that nothing else touches:
for an AND of $k+1$ variables a direct computation gives $E_i/\Inf_i=2(1-2^{-k})+2^k\Phi\bigl((1-2^{-k})^2,4^{-k}\bigr)$
at each of them (with $\Phi$ as in Section~\ref{sec:sep}), which differs from $2$ by at most $2^{-k}(2k+2)$. Theorem~\ref{thm:main} says that in a sparse
linear DNF every coordinate behaves like that, however many terms pass through it and however the terms around it
overlap.

The proof separates two spectra. Write $g$ and $h$ for the even and odd parts of $f$ in $x_i$. For a monotone $F$, $h$
is the indicator that $i$ is pivotal and $g=2N-1-h$, where $N$ is the indicator that $F$ vanishes at $x_i=0$. If the
spectrum of $h$ lives on a family $\Omega$ of sets on which $N$ has almost no weight, then on $\Omega$ the two
coefficients $\hat g(T)$ and $\hat h(T)$ have nearly the same square, and each such $T$ contributes nearly
$2\hat h(T)^2$ to $E_i$. Proposition~\ref{prop:sep} makes this quantitative for every monotone function. For a linear
DNF the family is $\Omega_k$, the sets that contain at least $k$ of the other variables of some term through $i$: the
pivotal indicator spreads its spectral weight evenly over the subsets of the other variables of each such term (Lemma~\ref{lem:h}), while the background
event $N$ needs $k$ further terms to reach $k$ of these variables, which costs $((D-1)2^{-(w-1)})^k$ (Lemma~\ref{lem:N}).

Theorem~\ref{thm:main} is a statement about the ratio at every coordinate and does not prove the FEI conjecture for the
class. The equivalence of~\cite{Shm1} needs a good coordinate in every restriction of $f$, and restrictions of a
uniform linear DNF have terms of different widths.

\section{A separation inequality}\label{sec:sep}

Fix $i$ and write $f^{(b)}$ for the restriction $x_i=b$, a function of the other coordinates. Put
\[
g=\tfrac12\bigl(f^{(0)}+f^{(1)}\bigr),\qquad h=\tfrac12\bigl(f^{(0)}-f^{(1)}\bigr).
\]
Then $\hat g(T)=\hat f(T)$ and $\hat h(T)=\hat f(T\cup\{i\})$ for $T\subseteq[n]\setminus\{i\}$, and $\E
h^2=\Inf_i(f)=:\delta$. Put $p(T)=\hat g(T)^2$ and $q(T)=\hat h(T)^2$, and
\[
\Phi(a,b)=\ell(a)+\ell(b)-\ell(a+b),\qquad \ell(x)=-x\log_2x.
\]
Since $(U\setminus\{i\},\mathbf 1[i\in U])$ takes the value $(T,0)$ with probability $p(T)$ and $(T,1)$ with probability
$q(T)$, we have $E_i(f)=\sum_T\Phi(p(T),q(T))$. The function $\Phi$ is concave, homogeneous of degree one, increasing in
each argument, and $\Phi(b,b)=2b$. Writing $\Phi(a,b)=a\log_2(1+b/a)+b\log_2(1+a/b)$ and using $\ln(1+t)\le t$ gives
\begin{equation}\label{eq:phiq}
\Phi(a,b)\le b\log_2\Bigl(e\,\frac{a+b}{b}\Bigr).
\end{equation}

Now let $F$ be monotone increasing and $f=1-2F$. Put $N=(1+f^{(0)})/2$ and $M=(1+f^{(1)})/2$, the indicators that $F$
vanishes at $x_i=0$ and at $x_i=1$, and $\Var N=\sum_{T\ne\emptyset}\hat N(T)^2$ is the variance of $N$. Then $M\le N$, $h=N-M$ is the indicator that $i$ is pivotal, $\hat
h(\emptyset)=\E h=\delta$, and
\[
g=N+M-1=2N-1-h,\qquad \hat g(T)=2\hat N(T)-\hat h(T)\quad(T\neq\emptyset).
\]

\begin{proposition}\label{prop:sep}
Let $F$ be monotone, $f=1-2F$, $\delta=\Inf_i(f)>0$, and let $N$ be as above. Let $\Omega$ be a family of non-empty
subsets of $[n]\setminus\{i\}$ and let $\varepsilon\in(0,1]$ and $\kappa\ge0$ satisfy
\[
\sum_{T\notin\Omega}\hat h(T)^2\le\varepsilon\delta,\qquad \sum_{T\in\Omega}\hat N(T)^2\le\kappa\delta.
\]
Then
\[
2-2\varepsilon-\frac{8}{\ln2}\sqrt\kappa\;\le\;\frac{E_i(f)}{\delta}\;\le\;2+\frac{8}{\ln2}\sqrt\kappa
+\varepsilon\log_2\Bigl(e\Bigl[3+\frac{8\Var N}{\varepsilon\delta}\Bigr]\Bigr)+\delta\log_2\bigl(e[1+\delta^{-2}]\bigr).
\]
\end{proposition}

\begin{proof}
Split the sum $E_i=\sum_T\Phi(p(T),q(T))$ into $T\in\Omega$, $T\notin\Omega$ with $T\neq\emptyset$, and $T=\emptyset$.

On $\Omega$. Fix $b\ge0$. The function $u\mapsto\Phi(u^2,b^2)$ is even, equals $2b^2$ at $u=\pm b$, and for $u\neq0$ has
derivative $2u\log_2(1+b^2/u^2)$, whose absolute value is at most $(4/\ln2)b$ because $\ln(1+t^2)\le2t$ for $t\ge0$.
For $T\in\Omega$ take $b=|\hat h(T)|$ and move $u$ from $-\hat h(T)$ to $\hat g(T)=2\hat N(T)-\hat h(T)$, a distance
$2|\hat N(T)|$:
\[
\bigl|\Phi(p(T),q(T))-2q(T)\bigr|\le\frac{8}{\ln2}|\hat h(T)||\hat N(T)|.
\]
By Cauchy--Schwarz, $\sum_\Omega|\hat h||\hat N|\le\sqrt\delta\sqrt{\kappa\delta}$, and $(1-\varepsilon)\delta\le
\sum_\Omega q\le\delta$. Hence
\[
2(1-\varepsilon)\delta-\frac{8}{\ln2}\delta\sqrt\kappa\le\sum_{T\in\Omega}\Phi(p(T),q(T))\le
2\delta+\frac{8}{\ln2}\delta\sqrt\kappa.
\]
Since every term of $E_i$ is non-negative, this gives the lower bound.

Off $\Omega$. Let $P$ and $Q$ be the sums of $p$ and $q$ over the non-empty $T\notin\Omega$; if $Q=0$ these terms
vanish, so let $Q>0$. Then $Q\le\varepsilon\delta$,
and from $p=(2\hat N-\hat h)^2\le8\hat N^2+2\hat h^2$ we get $P\le8\Var N+2\varepsilon\delta$. By~\eqref{eq:phiq} and
Jensen's inequality for the concave function $t\mapsto\log_2t$,
\[
\sum\Phi(p,q)\le\sum q\log_2\Bigl(e\,\frac{p+q}{q}\Bigr)\le Q\log_2\Bigl(e\,\frac{P+Q}{Q}\Bigr)\le
Q\log_2\Bigl(e\,\frac{c}{Q}\Bigr),\qquad c=8\Var N+3\varepsilon\delta.
\]
The last expression increases in $Q$ on $(0,c]$, so it is at most
$\varepsilon\delta\log_2(e[3+8\Var N/(\varepsilon\delta)])$.

At $T=\emptyset$, $q=\delta^2$ and $p\le1$, so~\eqref{eq:phiq} gives $\Phi\le\delta^2\log_2(e[1+\delta^{-2}])$. Adding
the three parts gives the upper bound.
\end{proof}

\section{Linear DNFs}\label{sec:lin}

From now on $F$ is a $w$-uniform linear monotone DNF of degree at most $D$, $f=1-2F$, $L=w-1\ge1$ and $a=2^{-L}$. Fix
a coordinate $i$ that lies in $d\ge1$ terms $t_1,\dots,t_d$. The \emph{mates} of $t_s$ are its $L$ variables other than
$i$, and $A_s$ is the indicator that all of them equal $1$. By linearity the mate sets are pairwise disjoint. A
\emph{background term} is a term that does not contain $i$; $N$ is the indicator that every background term fails, and
\[
M=N\prod_{s=1}^d(1-A_s),\qquad h=N-M=N\cdot\mathbf 1\bigl[\textstyle\bigvee_sA_s\bigr],\qquad \delta=\E h.
\]
Put $A=\E N$.

Two facts about linear DNFs are used repeatedly. A background term meets the mate set of $t_s$ in at most one variable,
since it shares at most one variable with $t_s$ and does not contain $i$. And a term through a mate of $t_s$, other than
$t_s$ itself, is a background term, since otherwise it would share two variables with $t_s$. We also use Harris's
inequality~\cite{Harris}: under a product measure, an increasing and a decreasing event are negatively correlated, and
two increasing events are positively correlated.

\begin{lemma}\label{lem:A}
With $\theta=1-L(D-1)a-(d-1)a/2$,
\[
da\,\theta\,A\le\delta\le da\,A.
\]
\end{lemma}

\begin{proof}
The events $A_s$ are increasing and $N$ is decreasing, so $\Pr[N,A_s]\le Aa$, and the union bound gives $\delta\le
daA$. For the lower bound, let $R_s$ be the indicator that every background term disjoint from the mate set of $t_s$
fails, so that $N\le R_s$. The remaining background terms meet that mate set in exactly one variable; each mate lies in
at most $D-1$ of them, so there are at most $L(D-1)$, and given $A_s=1$ each of them fires exactly when its other $L$
variables, which are not mates of $t_s$, all equal $1$. Let $V_s$ be the event that one of these residual conjunctions
equals $1$; it is increasing and $\Pr[V_s]\le L(D-1)a$. Since $R_s$ and $V_s$ do not depend on the mates of $t_s$,
\[
\Pr[N,A_s]=a\Pr[R_s,\overline{V_s}]\ge a\bigl(\Pr[R_s]-\Pr[R_s]\Pr[V_s]\bigr)\ge aA\bigl(1-L(D-1)a\bigr),
\]
using Harris for $R_s$ and $V_s$ and, in the last step, $\Pr[R_s]\ge A$; if $1-L(D-1)a<0$ the bound holds trivially.
For $s\neq u$ the event $A_sA_u$ is increasing with probability $a^2$, so $\Pr[N,A_s,
A_u]\le Aa^2$. Bonferroni's inequality gives $\delta\ge daA(1-L(D-1)a)-\binom d2a^2A=da\theta A$.
\end{proof}

For $1\le k\le L$ let $\Omega_k$ be the family of sets $T\subseteq[n]\setminus\{i\}$ that contain at least $k$ mates of
some $t_s$, and put
\[
\gamma_k=a\sum_{j<k}\binom Lj.
\]

\begin{lemma}\label{lem:h}
$\sum_{T\notin\Omega_k}\hat h(T)^2\le\min(1,d\gamma_k)\,\delta$.
\end{lemma}

\begin{proof}
Put $h_s=N A_s\prod_{u<s}(1-A_u)$. These functions have disjoint supports and sum to $h$, so $\sum_s\|h_s\|_2^2=\delta$.
Since $h_s$ vanishes unless $A_s=1$, $h_s=A_s\Psi_s$, where $\Psi_s$ is $N\prod_{u<s}(1-A_u)$ with the mates of $t_s$
set to $1$, a function of the variables outside that mate set. As
$A_s=2^{-L}\sum_{S\subseteq\text{mates}(t_s)}(-1)^{|S|}\chi_S$, every Fourier coefficient of $h_s$ has the form $\hat
h_s(S\cup V)=\pm2^{-L}\hat\Psi_s(V)$ with $S$ a set of mates and $V$ disjoint from them. The weight of $\hat h_s^2$ on
sets with exactly $j$ mates of $t_s$ is therefore $\binom Lj4^{-L}\|\Psi_s\|_2^2=\binom Lja\|h_s\|_2^2$, using
$\|h_s\|_2^2=a\|\Psi_s\|_2^2$. A set outside $\Omega_k$ has fewer than $k$ mates of $t_s$, so
$\sum_{T\notin\Omega_k}\hat h_s(T)^2\le\gamma_k\|h_s\|_2^2$. By Cauchy--Schwarz,
\[
\sum_{T\notin\Omega_k}\hat h(T)^2\le d\sum_s\sum_{T\notin\Omega_k}\hat h_s(T)^2\le d\gamma_k\delta.
\]
The bound by $\delta$ is trivial.
\end{proof}

\begin{lemma}\label{lem:N}
If $k(D-1)a<2$, then
\[
\sum_{T\in\Omega_k}\hat N(T)^2\le d\,\alpha_kA,\qquad
\alpha_k=\binom Lk\frac{2^{\binom k2-2k}\bigl((D-1)a\bigr)^k}{1-k(D-1)a/2}.
\]
\end{lemma}

\begin{proof}
Every $T\in\Omega_k$ contains a set $J$ of $k$ mates of some $t_s$, and there are $d\binom Lk$ such sets $J$. For each
of them, $\sum_{T\supseteq J}\hat N(T)^2=\|D_JN\|_2^2$, where $D_J$ is the composition of the derivatives
$D_jN=\frac12(N|_{x_j=0}-N|_{x_j=1})$ over $j\in J$. So it suffices to show $\|D_JN\|_2^2\le\alpha_k\binom Lk^{-1}A$.

Fix $J$. A background term contains at most one element of $J$. Let $R'$ be the indicator that every background term
avoiding $J$ fails, and for $j\in J$ let $Q_j$ be the OR of the residual conjunctions, of width $L$, of the at most
$D-1$ background terms through $j$. Neither $R'$ nor any $Q_j$ depends on the variables in $J$, and
\[
N=R'\prod_{j\in J}(1-x_jQ_j),\qquad D_JN=\pm R'\prod_{j\in J}\tfrac12Q_j,\qquad \|D_JN\|_2^2=4^{-k}\Pr[R',\textstyle\bigwedge_jQ_j].
\]
The event $\bigwedge_jQ_j$ is increasing and $R'$ decreasing, so $\Pr[R',\bigwedge_jQ_j]\le\Pr[R']\Pr[\bigwedge_jQ_j]$.
For $\bigwedge_jQ_j$ one residual conjunction must equal $1$ for each $j$; there are at most $(D-1)^k$ choices, the
chosen terms are distinct because each contains exactly one element of $J$, and two residual conjunctions share at most
one variable, so their union has at least $kL-\binom k2$ variables. Hence $\Pr[\bigwedge_jQ_j]\le(D-1)^k2^{\binom
k2}a^k$. Finally, $x_j$ is independent of $R'$ and $Q_j$, so by Harris
\[
A\ge\Pr[R']-\sum_{j\in J}\Pr[R',x_jQ_j]\ge\Pr[R']\Bigl(1-\sum_j\tfrac12\Pr[Q_j]\Bigr)\ge\Pr[R']\bigl(1-\tfrac12k(D-1)a\bigr).
\]
Combining the three displays gives $\|D_JN\|_2^2\le4^{-k}(D-1)^k2^{\binom k2}a^kA/(1-k(D-1)a/2)$.
\end{proof}

\begin{theorem}\label{thm:lin}
Let $F$ be a $w$-uniform linear monotone DNF of degree at most $D$, $L=w-1$, $a=2^{-L}$, and let $i$ lie in $d\ge1$
terms. Let $1\le k\le L$, and put
\[
\theta=1-L(D-1)a-\frac{(d-1)a}2,\qquad \varepsilon=\min(1,d\gamma_k),\qquad \kappa=\frac{\alpha_k}{a\theta},
\]
with $\gamma_k$ and $\alpha_k$ as above. If $\theta>0$, $k(D-1)a<2$ and $da\le1/4$, then
\[
\Bigl|\frac{E_i(f)}{\Inf_i(f)}-2\Bigr|\le2\varepsilon+\frac8{\ln2}\sqrt\kappa
+\varepsilon\log_2\Bigl(e\Bigl[3+\frac{8}{\varepsilon\,da\,\theta}\Bigr]\Bigr)+da\log_2\bigl(e[1+(da)^{-2}]\bigr).
\]
\end{theorem}

\begin{proof}
Apply Proposition~\ref{prop:sep} with $\Omega=\Omega_k$, which consists of non-empty sets. Lemma~\ref{lem:h} gives the
first hypothesis with $\varepsilon$ (note $\varepsilon\ge da>0$). Lemmas~\ref{lem:N} and~\ref{lem:A} give
$\sum_{\Omega_k}\hat N^2\le d\alpha_kA\le\alpha_k\delta/(a\theta)=\kappa\delta$. In the upper bound, $\Var N\le\E N=A\le
\delta/(da\theta)$, and $\delta\le daA\le da\le1/4$; the function $x\mapsto x\log_2(e[1+x^{-2}])$ is increasing on
$(0,1/4]$, since its derivative is $\log_2(1+x^{-2})+\log_2e\,(1-2/(1+x^2))\ge\log_2 17-\log_2e>0$ there.
\end{proof}

\begin{proof}[Proof of Theorem~\ref{thm:main}]
Let $k=\lfloor1/\eta\rfloor+1$, so $k\eta>1$, and let $L\to\infty$ with $D\le2^{(1-\eta)L}$. Then $Da\le2^{-\eta L}$,
so $L(D-1)a\to0$, $\theta\to1$, $k(D-1)a<2$ and $da\le Da\le1/4$ for $L$ large, uniformly in $d\le D$. Next
$d\gamma_k\le Da\,kL^k\le kL^k2^{-\eta L}\to0$, and
\[
\kappa\le\frac{2L^k2^{\binom k2}(Da)^k}{a\theta}\le\frac{2L^k2^{\binom k2}}{\theta}\,2^{(1-k\eta)L}\to0.
\]
In the third term, $x\mapsto x\log_2(e[3+c/x])$ is increasing, $\varepsilon\le kL^k2^{-\eta L}$, and $c=8/(da\theta)\le
16\cdot2^L$ for $L$ large, so the term is $O(kL^{k+1}2^{-\eta L})$. The last term is at most
$2^{-\eta L}\log_2(e[1+4^{\eta L}])=O(L2^{-\eta L})$. None of these bounds depends on $n$, on the number of terms or on
$i$.
\end{proof}

\begin{remark}
The two error sources pull $k$ in opposite directions. Since $\gamma_k\le kL^{k-1}a$, the weight of $h$ off $\Omega_k$ is
small once $DkL^{k-1}\ll2^L$, which favours small $k$; the weight of $N$ on $\Omega_k$ is small once
$L^kD^k\ll2^{(k-1)L}$ up to factors depending on $k$, which favours large $k$. With $D=2^{\beta L}$, $\beta<1$, both
hold for $k>1/(1-\beta)$, and this is where the hypothesis $D\le2^{(1-\eta)L}$ comes from.
\end{remark}

\section{Numerical checks}\label{sec:num}

All values below were computed with exact Walsh--Hadamard transforms. The script \texttt{check\_lin.py} in the
code repository evaluates the four inequalities of Lemmas~\ref{lem:A}--\ref{lem:N} and Theorem~\ref{thm:lin} at every coordinate and every
admissible $k$ for $17$ linear DNFs: stars of $D$ terms through one variable, the rows and columns of a $4\times4$
array, and twelve random linear DNFs with $w\in\{4,5\}$, $D\in\{2,3,4\}$ and $16$ to $18$ variables. All $784$
(coordinate, $k$) pairs satisfy all four inequalities, and the left side of Theorem~\ref{thm:lin} is never more than
$0.14$ of the right side; at these widths the bound is far from tight.

Table~\ref{tab:star} shows the ratio approaching $2$ as the width grows, for the star of $D$ terms of width $w$ through
a common variable, at the centre ($d=D$) and at another variable ($d=1$). For the rows and columns of an $m\times m$
array the ratio is $2.629$ at $m=3$ and $3.245$ at $m=4$.

\begin{table}[tbp]
\centering
\begin{tabular}{rrrcccc}
\toprule
$D$ & $w$ & $n$ & $\Inf$ centre & $E/\Inf$ centre & $\Inf$ other & $E/\Inf$ other\\\midrule
2 & 4 & 7 & 0.2344 & 2.685 & 0.1094 & 2.909\\
2 & 6 & 11 & 0.0615 & 2.450 & 0.0303 & 2.351\\
2 & 8 & 15 & 0.0156 & 2.178 & 0.0078 & 2.118\\
2 & 10 & 19 & 0.0039 & 2.060 & 0.0019 & 2.037\\
3 & 5 & 13 & 0.1760 & 2.692 & 0.0549 & 2.666\\
3 & 7 & 19 & 0.0461 & 2.378 & 0.0151 & 2.225\\
4 & 5 & 17 & 0.2275 & 2.690 & 0.0515 & 2.728\\
\bottomrule
\end{tabular}
\caption{$E_i/\Inf_i$ for the star of $D$ terms of width $w$ on $n=1+D(w-1)$ variables.}\label{tab:star}
\end{table}

\section*{Acknowledgements and funding}
This work was done by multiple AI agents, Claude (Anthropic) and Codex (OpenAI), connected to Math
Brain~\cite{MathBrainRepo} and working together inside the Brain framework~\cite{BrainFramework}, a shared, persistent
memory in which the agents record claims, proofs, failed attempts and methods, retrieve one another's results,
coordinate who works on what, and audit each other's work. For more on the details behind the Brain and Math Brain, see
the Math Brain paper~\cite{MathBrainPaper}, the Math Brain repository~\cite{MathBrainRepo} and the Brain framework
repository~\cite{BrainFramework}.

The research presented in this paper was supported by the European Research Council (ERC) under the European Union's
Horizon Europe research and innovation programme (grant agreement No.~101041711), by the Simons Foundation as part of
the Collaboration on the Mathematical and Scientific Foundations of Deep Learning, by Heights Labs, by Convex Nexus
Capital, by the Israel Science Foundation (grant number 2258/19), and by the Israel Science Foundation (ISF Grant
4101/25). Computational resources were provided in part by Google Cloud through its education credits programme (credit
530550103).

\section*{Declaration of competing interest}
The author holds equity in Heights Labs and in Convex Nexus Capital and has a management role at Convex Nexus Capital.
These entities are included in the funding declaration above.

\section*{Declaration of generative AI and AI-assisted technologies}
The paper comes out of Math Brain, a project on the use of AI systems in research mathematics~\cite{BrainFramework,MathBrainRepo}. In November 2026 a
Claude Fable agent (Anthropic) and Codex agents running OpenAI's GPT-6.1, connected to the Brain and working on the same
problem, studied erasure entropies of DNFs; a Codex agent found the theorem for sparse linear DNFs and its proof.
Claude Opus 5.5 (Anthropic) rederived each proof, wrote the numerical checks, prepared the manuscript and audited it.
The author directed the work and is responsible for the content of the paper.

\section*{Code and data availability}
Code and outputs: \url{https://github.com/yspennstate/fei-sparse-linear-dnf}, archived on Zenodo with doi:10.5281/zenodo.24125386.

The paper is part of the Math Brain project, whose public repository is
\url{https://github.com/yspennstate/math-brain-public}~\cite{MathBrainRepo}.

\appendix
\section{Notation}\label{app:notation}

Symbols are listed roughly in the order in which they first appear, and the last column gives that place. The letters
$c$, $k$ and $S$ are reused, and each meaning has its own row; $D$ is the degree bound, and $D_j$, $D_J$ are derivatives.

\begin{longtable}{@{}>{\raggedright\arraybackslash}p{0.19\textwidth}>{\raggedright\arraybackslash}p{0.54\textwidth}>{\raggedright\arraybackslash}p{0.2\textwidth}@{}}
\toprule
Symbol & Meaning & First appears \\
\midrule
\endhead
$f$ & a Boolean function $\{0,1\}^n\to\{-1,1\}$; for a DNF $F$, $f=1-2F$ & Section~\ref{sec:intro} \\
$\hat f(S)$, $\chi_S$ & the Fourier coefficient of $f$ at $S$, and the character $(-1)^{\sum_{j\in S}x_j}$ & Section~\ref{sec:intro} \\
$U$ & a random subset of $[n]$ with $\Pr[U=S]=\hat f(S)^2$ & Section~\ref{sec:intro} \\
$H(f)$ & the Fourier entropy $H(U)$ & Section~\ref{sec:intro} \\
$\Inf_i(f)$, $I(f)$ & the influence $\Pr[i\in U]$ of $i$, and the total influence $\E|U|$ & Section~\ref{sec:intro} \\
$C$ & an absolute constant, as in $H(f)\le C\,I(f)$ & Section~\ref{sec:intro} \\
$k$ & the number of times a decision tree or DNF may read a variable, in read-$k$ & Section~\ref{sec:intro} \\
$E_i(f)$ & the erasure entropy $H(\mathbf 1[i\in U]\mid U\setminus\{i\})$ & Section~\ref{sec:intro} \\
$c$ & an absolute constant, as in $E_i(f)\le c\,\Inf_i(f)$ & Section~\ref{sec:intro} \\
$h_2$ & the binary entropy & Section~\ref{sec:intro} \\
$F$ & a monotone DNF on $\{0,1\}^n$ & Section~\ref{sec:intro} \\
$w$, $D$ & the width of every term, and the largest number of terms through one variable & Section~\ref{sec:intro} \\
$m$ & the side of the $m\times m$ array whose rows and columns are the terms & Section~\ref{sec:intro} \\
$\eta$ & a fixed number in $(0,1)$, as in $D\le2^{(1-\eta)(w-1)}$ & Theorem~\ref{thm:main} \\
$k$ & the number of variables of an AND, less one, in the comparison with a single AND & Section~\ref{sec:intro} \\
$\Phi(a,b)$ & $\ell(a)+\ell(b)-\ell(a+b)$, with $\ell$ as below & Section~\ref{sec:intro} \\
$g$, $h$ & the even and odd parts $\frac12(f^{(0)}\pm f^{(1)})$ of $f$ in $x_i$ & Section~\ref{sec:intro} \\
$N$ & the indicator that $F$ vanishes at $x_i=0$ (for a linear DNF, that every background term fails) & Section~\ref{sec:intro} \\
$\Var N$ & the variance $\sum_{T\ne\emptyset}\hat N(T)^2$ of $N$ & Section~\ref{sec:sep} \\
$\Omega$ & a family of non-empty sets on which $h$ has most of its weight and $N$ little & Section~\ref{sec:intro} \\
$\Omega_k$ & the sets containing at least $k$ mates of some term through $i$ & Section~\ref{sec:intro} \\
$f^{(b)}$ & the restriction $x_i=b$ & Section~\ref{sec:sep} \\
$T$ & a subset of $[n]\setminus\{i\}$ & Section~\ref{sec:intro} \\
$\delta$ & the influence $\Inf_i(f)=\E h^2$ & Section~\ref{sec:sep} \\
$p(T)$, $q(T)$ & $\hat g(T)^2$ and $\hat h(T)^2$ & Section~\ref{sec:sep} \\
$\ell(x)$ & $-x\log_2x$ & Section~\ref{sec:sep} \\
$M$ & the indicator that $F$ vanishes at $x_i=1$ & Section~\ref{sec:sep} \\
$\varepsilon$, $\kappa$ & the weight of $h$ off $\Omega$ and the weight of $N$ on $\Omega$, as fractions of $\delta$ & Proposition~\ref{prop:sep} \\
$P$, $Q$ & the sums of $p$ and $q$ over the non-empty sets outside $\Omega$, in the proof & Proposition~\ref{prop:sep} \\
$c$ & $8\Var N+3\varepsilon\delta$, in the proof & Proposition~\ref{prop:sep} \\
$L$, $a$ & $w-1$ and $2^{-L}$ & Section~\ref{sec:lin} \\
$d$, $t_1,\dots,t_d$ & the number of terms through $i$, and those terms & Section~\ref{sec:lin} \\
mates, $A_s$ & the $L$ variables of $t_s$ other than $i$, and the indicator that all of them equal $1$ & Section~\ref{sec:lin} \\
background term & a term that does not contain $i$ & Section~\ref{sec:lin} \\
$A$ & $\E N$ & Section~\ref{sec:lin} \\
$\theta$ & $1-L(D-1)a-(d-1)a/2$ & Lemma~\ref{lem:A} \\
$R_s$, $V_s$ & the indicator that every background term disjoint from the mates of $t_s$ fails, and the event that one residual conjunction of the others equals $1$ & Lemma~\ref{lem:A} \\
$k$ & the number of mates that puts a set in $\Omega_k$, $1\le k\le L$ & Section~\ref{sec:lin} \\
$\gamma_k$ & $a\sum_{j<k}\binom Lj$ & Section~\ref{sec:lin} \\
$\beta$ & the exponent in $D=2^{\beta L}$, $\beta<1$, in the closing remark & Section~\ref{sec:lin} \\
$h_s$, $\Psi_s$ & $NA_s\prod_{u<s}(1-A_u)$, and the same product with the mates of $t_s$ set to $1$ & Lemma~\ref{lem:h} \\
$S$, $V$ & a set of mates of $t_s$, and a set disjoint from them, in the proof & Lemma~\ref{lem:h} \\
$\alpha_k$ & the constant bounding the weight of $N$ on $\Omega_k$ & Lemma~\ref{lem:N} \\
$J$ & a set of $k$ mates of one term through $i$ & Lemma~\ref{lem:N} \\
$D_j$, $D_J$ & the derivative $\frac12(N|_{x_j=0}-N|_{x_j=1})$ in $x_j$, and the composition of these over $j\in J$ & Lemma~\ref{lem:N} \\
$R'$, $Q_j$ & the indicator that every background term avoiding $J$ fails, and the OR of the residual conjunctions of the background terms through $j$ & Lemma~\ref{lem:N} \\
\bottomrule
\end{longtable}

\begin{thebibliography}{99}
\bibitem{ACKSW} S.~Arunachalam, S.~Chakraborty, M.~Kouck\'y, N.~Saurabh, R.~de~Wolf, Improved bounds on Fourier
entropy and min-entropy, ACM Trans. Comput. Theory 13 (2021), no.~4, 1--40, doi:10.1145/3470860; arXiv:1809.09819.
\bibitem{CKLS} S.~Chakraborty, R.~Kulkarni, S.~V.~Lokam, N.~Saurabh, Upper bounds on Fourier entropy, Theoret.
Comput. Sci. 654 (2016), 92--112.
\bibitem{FK96} E.~Friedgut, G.~Kalai, Every monotone graph property has a sharp threshold, Proc. Amer. Math. Soc.
124 (1996), 2993--3002.
\bibitem{GMP} M.~J.~Gonz\'alez, P.~MacManus, M.~C.~Pereyra, Further evidence towards the Fourier
entropy-influence conjecture, arXiv:2606.00246, 2026.
\bibitem{Han} X.~Han, A new bound for the Fourier-entropy-influence conjecture, Combinatorica 45 (2025), no.~1,
Art.~4, doi:10.1007/s00493-024-00133-z; arXiv:2312.08271.
\bibitem{Harris} T.~E.~Harris, A lower bound for the critical probability in a certain percolation process, Proc.
Cambridge Philos. Soc. 56 (1960), 13--20.
\bibitem{KKLMS} E.~Kelman, G.~Kindler, N.~Lifshitz, D.~Minzer, M.~Safra, Towards a proof of the
Fourier-entropy conjecture?, Geom. Funct. Anal. 30 (2020), 1097--1138.
\bibitem{LH25} P.~Li, G.~Han, Strengthening Han's Fourier entropy-influence inequality via an
information-theoretic proof, arXiv:2512.03117, 2025.
\bibitem{OT13} R.~O'Donnell, L.-Y.~Tan, A composition theorem for the Fourier entropy-influence conjecture, ICALP 2013,
Lecture Notes in Comput. Sci. 7965, 780--791, doi:10.1007/978-3-642-39206-1\_66.
\bibitem{OWZ} R.~O'Donnell, J.~Wright, Y.~Zhou, The Fourier entropy-influence conjecture for certain classes of
Boolean functions, ICALP 2011, Lecture Notes in Comput. Sci. 6755, 330--341.
\bibitem{Shalev} G.~Shalev, On the Fourier entropy influence conjecture for extremal classes, arXiv:1806.03646, 2018.
\bibitem{Shm1} Y.~Shmalo, The one-coordinate step in the Fourier entropy-influence conjecture,
\url{https://github.com/yspennstate/fourier-entropy-influence-one-coordinate}, doi:10.5281/zenodo.24125404, 2026.
\bibitem{ShmLTF} Y.~Shmalo, Fourier entropy of threshold functions and their compositions,
arXiv:2610.16069, 2026.
\bibitem{VW08} S.~Verd\'u, T.~Weissman, The information lost in erasures, IEEE Trans. Inform. Theory 54 (2008),
5030--5058.
\bibitem{WWW} A.~Wan, J.~Wright, C.~Wu, Decision trees, protocols, and the Fourier entropy-influence conjecture, ITCS 2014,
arXiv:1312.3003.
\bibitem{MathBrainRepo} Y.~Shmalo, \emph{Math Brain: public snapshot, code and tests}, 2026,
\url{https://github.com/yspennstate/math-brain-public}.
\bibitem{MathBrainPaper} Y.~Shmalo, \emph{Math Brain: shared, auditable memory for multi-agent mathematical research},
October 2026, \url{https://github.com/yspennstate/math-brain-public/tree/main/papers/math-brain-2026-10-24}, archived on Zenodo with doi:10.5281/zenodo.24150837.
\bibitem{BrainFramework} Y.~Shmalo, \emph{Brain Framework}, public software repository, 2026,
\url{https://github.com/yspennstate/brain-framework-public}.
\end{thebibliography}

\end{document}
